% ═══════════════════════════════════════════════════════════════
% MUSTERLÖSUNG: Radioaktivität (Physik Klasse 10)
% Identisches Layout wie AB, Lösungen in ROT
% ═══════════════════════════════════════════════════════════════

\documentclass[10pt, a4paper]{article}

\usepackage[utf8]{inputenc}
\usepackage[T1]{fontenc}
\usepackage[ngerman]{babel}
\usepackage{helvet}
\renewcommand{\familydefault}{\sfdefault}

\usepackage[margin=1.5cm, top=1.8cm, bottom=1.5cm]{geometry}
\usepackage{enumitem}
\usepackage{tabularx}
\usepackage{booktabs}
\usepackage{array}
\usepackage{multicol}

\usepackage{xcolor}
\usepackage{tcolorbox}
\usepackage{amssymb}
\usepackage{fancyhdr}
\usepackage{lastpage}

\setlist{nosep, leftmargin=1.5em}

% Farben
\definecolor{physik}{HTML}{2c5aa0}
\definecolor{grau}{HTML}{888888}
\definecolor{hellgrau}{HTML}{f0f0f0}
\definecolor{fehler}{HTML}{c62828}
\definecolor{fehlerbg}{HTML}{ffebee}
\definecolor{loesung}{HTML}{c62828}

\newcommand{\fachfarbe}{physik}

% Variablen
\newcommand{\thema}{Radioaktivität}
\newcommand{\fach}{Physik}
\newcommand{\klasse}{Klasse 10}
\newcommand{\stunde}{Stunde 3-4}

% Kopf- und Fußzeile
\pagestyle{fancy}
\fancyhf{}
\renewcommand{\headrulewidth}{0.4pt}
\lhead{\small\fach{} \klasse{} -- \thema{} \textcolor{loesung}{\textbf{[MUSTERLÖSUNG]}}}
\rhead{\small Name: \rule{4cm}{0.4pt}}
\cfoot{\small\thepage/\pageref{LastPage}}

% Befehle
\newcommand{\lsg}[1]{\textcolor{loesung}{\textbf{#1}}}

\newtcolorbox{merksatz}[1][]{
    colback=hellgrau, colframe=\fachfarbe, boxrule=0.8pt, arc=0pt,
    left=6pt, right=6pt, top=6pt, bottom=6pt,
    fonttitle=\bfseries\small, title=#1
}

\newtcolorbox{fehlerbox}{
    colback=fehlerbg, colframe=fehler, boxrule=0.8pt, arc=0pt,
    left=6pt, right=6pt, top=4pt, bottom=4pt,
    fonttitle=\bfseries\small, title=Typische Fehler
}

\newcommand{\abschnitt}[1]{%
    \vspace{8pt}\noindent\textbf{\textcolor{\fachfarbe}{#1}}\vspace{4pt}
}

\begin{document}

\begin{center}
{\large\bfseries Arbeitsblatt: \thema{} -- \stunde{} \textcolor{loesung}{[MUSTERLÖSUNG]}}\\[2pt]
{\small\textcolor{grau}{Begleitend zum Lernpfad -- in den Hefter übernehmen}}
\end{center}
\vspace{-2pt}\hrule\vspace{8pt}

\begin{multicols}{2}
\small

% ─────────────────────────────────────────────────────────────
% KASTEN 1: Die drei Strahlungsarten
% ─────────────────────────────────────────────────────────────

\abschnitt{Kasten 1: Die drei Strahlungsarten}

\begin{merksatz}[Vergleichstabelle]
Ergänze die Tabelle aus dem Lernpfad:

\vspace{4pt}
\renewcommand{\arraystretch}{1.6}
\begin{tabular}{|p{1.8cm}|p{1.5cm}|p{1.5cm}|p{1.5cm}|}
\hline
\textbf{Eigenschaft} & \textbf{α} & \textbf{β} & \textbf{γ} \\
\hline
Natur & \lsg{He-Kerne} & \lsg{Elektronen} & \lsg{EM-Wellen} \\
\hline
Ladung & \lsg{+2} & \lsg{-1} & \lsg{0} \\
\hline
Reichweite (Luft) & \lsg{wenige cm} & \lsg{einige m} & \lsg{km} \\
\hline
Abschirmung & \lsg{Papier} & \lsg{Aluminium} & \lsg{Blei} \\
\hline
\end{tabular}
\renewcommand{\arraystretch}{1}
\end{merksatz}

\vspace{6pt}

% ─────────────────────────────────────────────────────────────
% KASTEN 2: Zerfallsgleichungen
% ─────────────────────────────────────────────────────────────

\abschnitt{Kasten 2: Zerfallsgleichungen}

\begin{merksatz}[Regeln für α- und β-Zerfall]
\textbf{α-Zerfall:}\\
$^{A}_{Z}$X $\rightarrow$ $^{\lsg{A-4}}_{\lsg{Z-2}}$Y + $^{4}_{2}$He

\vspace{4pt}
Regel: A sinkt um \lsg{4}, Z sinkt um \lsg{2}

\vspace{8pt}

\textbf{β$^-$-Zerfall:}\\
$^{A}_{Z}$X $\rightarrow$ $^{\lsg{A}}_{\lsg{Z+1}}$Y + $^{0}_{-1}$e

\vspace{4pt}
Regel: A bleibt \lsg{gleich}, Z steigt um \lsg{1}
\end{merksatz}

\vspace{6pt}

% ─────────────────────────────────────────────────────────────
% KASTEN 3: Ionisierende Wirkung
% ─────────────────────────────────────────────────────────────

\abschnitt{Kasten 3: Ionisierende Wirkung}

\begin{merksatz}[Definition]
Ionisierende Strahlung kann \lsg{Elektronen} aus

Atomen herausschlagen. Es entstehen \lsg{Ionen}.

\vspace{6pt}

Biologische Wirkung: \lsg{DNA-Schäden, Zellzerstörung}
\end{merksatz}

\columnbreak

% ─────────────────────────────────────────────────────────────
% ÜBUNG 1: Zuordnung (AFB I)
% ─────────────────────────────────────────────────────────────

\abschnitt{Übung 1: Eigenschaften zuordnen}

Ordne zu: α, β oder γ

\vspace{4pt}
\begin{tabular}{|p{4.5cm}|c|}
\hline
\textbf{Aussage} & \textbf{Strahlung} \\
\hline
Wird von Papier gestoppt & \lsg{α} \\
\hline
Besteht aus Heliumkernen & \lsg{α} \\
\hline
Elektromagnetische Wellen & \lsg{γ} \\
\hline
Höchstes Ionisierungsvermögen & \lsg{α} \\
\hline
Wird von Aluminium gestoppt & \lsg{β} \\
\hline
\end{tabular}

\vspace{8pt}

% ─────────────────────────────────────────────────────────────
% ÜBUNG 2: Zerfallsgleichungen (AFB II)
% ─────────────────────────────────────────────────────────────

\abschnitt{Übung 2: Zerfallsgleichungen aufstellen}

Stelle die Zerfallsgleichungen auf:

\vspace{6pt}

a) $^{238}_{92}$U zerfällt unter α-Strahlung:

\vspace{4pt}
\lsg{$^{238}_{92}$U $\rightarrow$ $^{234}_{90}$Th + $^{4}_{2}$He}

\vspace{6pt}

b) $^{14}_{6}$C zerfällt unter β$^-$-Strahlung:

\vspace{4pt}
\lsg{$^{14}_{6}$C $\rightarrow$ $^{14}_{7}$N + $^{0}_{-1}$e}

\vspace{6pt}

c) $^{226}_{88}$Ra zerfällt unter α-Strahlung:

\vspace{4pt}
\lsg{$^{226}_{88}$Ra $\rightarrow$ $^{222}_{86}$Rn + $^{4}_{2}$He}

\vspace{8pt}

% ─────────────────────────────────────────────────────────────
% ÜBUNG 3: Fehler erklären (AFB III)
% ─────────────────────────────────────────────────────────────

\abschnitt{Übung 3: Fehler erklären (AFB III)}

\begin{fehlerbox}
Erkläre, warum diese Aussage falsch ist:

\vspace{4pt}

\textit{„α-Strahlung ist ungefährlich, weil sie schon von Papier gestoppt wird."}

\vspace{4pt}
\lsg{Falsch! α-Strahlung hat das höchste Ionisierungsvermögen. Bei Inkorporation (Einatmen, Verschlucken) ist sie besonders gefährlich, da sie lokal sehr viel Energie abgibt.}
\end{fehlerbox}

\vspace{8pt}

% ─────────────────────────────────────────────────────────────
% ÜBUNG 4: Transfer (AFB III)
% ─────────────────────────────────────────────────────────────

\abschnitt{Übung 4: Anwendung}

Welche Schutzmaßnahmen brauchst du bei einem γ-Strahler? Begründe.

\vspace{4pt}
\lsg{Dicke Bleiabschirmung (mehrere cm), großer Abstand zur Quelle, kurze Aufenthaltszeit. γ-Strahlung durchdringt die meisten Materialien, nur dichte Materialien wie Blei oder Beton können sie effektiv abschirmen.}

\end{multicols}

% ═══════════════════════════════════════════════════════════════
% SEITE 2: VERSUCHSPROTOKOLL
% ═══════════════════════════════════════════════════════════════

\newpage

\begin{center}
{\large\bfseries Versuchsprotokoll: Abschirmung radioaktiver Strahlung \textcolor{loesung}{[MUSTERLÖSUNG]}}\\[2pt]
{\small\textcolor{grau}{Demonstration mit Geiger-Müller-Zählrohr}}
\end{center}
\vspace{-2pt}\hrule\vspace{12pt}

\abschnitt{Teil A: Versuchsaufbau}

\begin{merksatz}[Skizze des Aufbaus]
Zeichne den Versuchsaufbau: Strahlungsquelle → Absorber → Geiger-Müller-Zählrohr

\vspace{6pt}
\lsg{Skizze: Links Strahlungsquelle (Behälter mit radioaktivem Symbol) $\rightarrow$ Mitte: Absorberplatte (senkrecht, austauschbar) $\rightarrow$ Rechts: Geiger-Müller-Zählrohr (zylindrisch) $\rightarrow$ Zählgerät mit Display}
\vspace{40pt}
\end{merksatz}

\vspace{8pt}

\abschnitt{Teil B: Vorhersage (VOR dem Versuch)}

Welche Strahlung wird von welchem Material gestoppt? Kreuze an:

\vspace{6pt}

\begin{tabular}{|l|c|c|c|}
\hline
\textbf{Material} & \textbf{stoppt α} & \textbf{stoppt β} & \textbf{stoppt γ} \\
\hline
Papier (1 mm) & \lsg{$\boxtimes$} & $\square$ & $\square$ \\
\hline
Aluminium (2 mm) & \lsg{$\boxtimes$} & \lsg{$\boxtimes$} & $\square$ \\
\hline
Blei (5 mm) & \lsg{$\boxtimes$} & \lsg{$\boxtimes$} & \lsg{$\boxtimes$} \\
\hline
\end{tabular}

\vspace{12pt}

\abschnitt{Teil C: Beobachtung (WÄHREND des Versuchs)}

Notiere die Zählraten (Impulse pro Minute):

\vspace{6pt}

\begin{tabular}{|l|c|c|}
\hline
\textbf{Absorber} & \textbf{Zählrate (Imp/min)} & \textbf{Veränderung} \\
\hline
Ohne Absorber & \lsg{ca. 500} & -- \\
\hline
Papier (1 mm) & \lsg{ca. 350} & \lsg{sinkt um ca. 30\%} \\
\hline
Aluminium (2 mm) & \lsg{ca. 120} & \lsg{sinkt um ca. 75\%} \\
\hline
Blei (5 mm) & \lsg{ca. 20} & \lsg{sinkt um ca. 96\%} \\
\hline
\end{tabular}

\vspace{12pt}

\abschnitt{Teil D: Auswertung}

\textbf{1. Welche Strahlungsart wurde hauptsächlich verwendet?}

\vspace{4pt}
\lsg{Am wahrscheinlichsten β-Strahlung oder ein Mischpräparat. Reine α-Strahlung wäre schon bei Papier fast komplett absorbiert.}

\vspace{8pt}

\textbf{2. Stimmt deine Vorhersage mit der Beobachtung überein? Erkläre Abweichungen.}

\vspace{4pt}
\lsg{Weitgehend ja. Abweichungen: Die Absorption ist nie 100\% (Restuntergrund), und γ-Anteile können auch Blei durchdringen.}

\vspace{8pt}

\textbf{3. Warum reicht bei α-Strahlern schon Papier als Schutz, bei γ-Strahlern aber nicht?}

\vspace{4pt}
\lsg{α-Teilchen sind schwer (Heliumkerne) und stark geladen (+2), daher werden sie schnell abgebremst. γ-Strahlung ist ungeladen und durchdringt Materie fast ungehindert, nur sehr dichte Materialien wie Blei absorbieren sie.}

\end{document}
